
complexes de Koszul associés. Le complexe de Koszul associé à la suite @@M0@@ est isomorphe au complexe produit tensoriel @@M1@@. Comme c'est une résolution de @@M2@@ par des @@M3@@-Modules libres de type fini, les @@M4@@ sont les faisceaux d'homologie du complexe @@M5@@. On peut le considérer comme le produit tensoriel des complexes @@M6@@ et @@M7@@, et utiliser pour calculer son homologie les suites spectrales associées à un bicomplexe (EGA @@M8@@ (11.3.5)), qui ont comme aboutissement l'homologie du complexe simple associé.

On considère la suite spectrale telle que @@M9@@, avec les notations de EGA @@M10@@ (11.3.5), le premier indice étant relatif au complexe @@M11@@ et le second au complexe @@M12@@. Or le complexe @@M13@@ est le complexe de Koszul sur @@M14@@ associé à la suite de sections @@M15@@, où les @@M16@@ sont les images des @@M17@@ dans @@M18@@, et ce complexe est une résolution de @@M19@@ sur @@M20@@ car les immersions @@M21@@, et @@M22@@ sont régulières (cf. démonstration de 1.7 ii). Comme @@M23@@ est libre sur @@M24@@, le complexe @@M25@@ est une résolution de @@M26@@. Par conséquent la suite spectrale dégénère, et les @@M27@@ sont les faisceaux d'homologie du complexe @@M28@@. Or la différentielle de @@M29@@ s'exprime en termes des @@M30@@, qui annulent @@M31@@; par suite, la différentielle de @@M32@@ est nulle, et on obtient donc les isomorphismes:
@@M33@@

En particulier, on trouve pour @@M34@@:
@@M35@@
d'où pour tout @@M36@@:
@@M37@@

